\documentclass[11pt]{article}
\usepackage[T1]{fontenc}
\usepackage{lmodern,microtype,amsmath,amssymb,amsthm,mathtools,mathrsfs,tikz-cd}
\usepackage[margin=1in]{geometry}
\usepackage{enumitem,booktabs,xurl,hyperref}
\hypersetup{colorlinks=true,linkcolor=blue,citecolor=blue,urlcolor=blue,
 pdftitle={Finite proof algebras in range restriction categories},pdfauthor={Florian Lengyel}}
\newtheorem{theorem}{Theorem}[section]
\newtheorem{proposition}[theorem]{Proposition}
\newtheorem{lemma}[theorem]{Lemma}
\newtheorem{corollary}[theorem]{Corollary}
\theoremstyle{definition}
\newtheorem{definition}[theorem]{Definition}
\theoremstyle{remark}
\newtheorem{remark}[theorem]{Remark}
\newtheorem{problem}[theorem]{Problem}
\newcommand{\R}{\mathscr R}
\newcommand{\D}{\mathsf D}
\newcommand{\N}{\mathbb N}
\newcommand{\ran}{\operatorname{ran}}
\newcommand{\Tot}{\operatorname{Tot}}
\newcommand{\Set}{\mathbf{Set}}
\newcommand{\op}{\mathrm{op}}
\newcommand{\id}{\operatorname{id}}
\title{Finite proof algebras\\in range restriction categories}
\author{Florian Lengyel}
\date{29 October 2026}
\begin{document}
\maketitle
\begin{abstract}
Increasing domains of finite constructor stages determine an initial
algebra when the corresponding partial folds admit gluing with exactly
those domains. A universal object with partial retractions supplies the
typed constructor embeddings. For finitary inference polynomials, a
Turing morphism constructs the partial folds by a parameterized recursion
theorem; the finite-stage domain unions remain an existence hypothesis.
Under regularity, parameter-stable finite
recursion, and decidable equality on atoms, disjoint computational
domains determine a consistent classical deductive presentation. An
inseparable pair makes every consistent computational extension of this
presentation incomplete. The stated construction hypotheses are realized
by partial recursive maps. The range question for a selected presheaf
completion is examined in relation to Heller's example of a recursion
category without ranges.
\end{abstract}
\tableofcontents

\section{Computational domains and inseparability}\label{gen:background}

Finite proofs require both an object of codes and a recursion principle
for interpreting those codes. In a category of partial computations,
the finite stages of an interpretation can be partial arrows with
increasing domains. Their gluing gives the fold required by initiality.
The construction below applies this principle to typed inference rules
and then proves consistency by finite Boolean valuations.

We use restriction categories in the sense of \cite{CockettLack}.
Composition is written $gf=g\circ f$. A restriction category assigns
to each $f:X\to Y$ a domain idempotent $\bar f:X\to X$ satisfying
\[
 f\bar f=f,\qquad \bar f\bar g=\bar g\bar f,
 \qquad\overline{g\bar f}=\bar g\bar f,
 \qquad\bar h f=f\overline{hf}.
\]
The middle identities require $f,g$ to have the same source; the last
requires $h$ to have source $Y$. An arrow is total when $\bar f=1_X$.
A restriction idempotent $e=\bar e$ is called a domain. We assume zero
arrows with $\overline{0_{X,Y}}=0_X$ whenever zeros are used. Domains
are ordered by $e\le d$ if $de=e$, with meet $ed$ and bottom $0_X$.

For $f:X\to Y$ and a domain $q$ on $Y$, inverse image is
$f^*q=\overline{qf}$. A range operator assigns a domain $\ran(f)$ on
$Y$ such that, for every domain $e$ on $Y$ and arrow $g:Y\to Z$,
\[
 \ran(f)f=f,\qquad
 \ran(ef)=e\ran(f),\qquad
 \ran(g\ran(f))=\ran(gf).
\]
In particular, $\ran(f)\le e$ if and only if $ef=f$.
Where specified, the universal range property further says that
$vf=wf$ implies $v\ran(f)=w\ran(f)$.

Write $X\boxtimes Y$ for a binary restriction product. Its projections
are total, and pairing satisfies
\[
 \pi_1\langle f,g\rangle=f\bar g,\qquad
 \pi_2\langle f,g\rangle=g\bar f,\qquad
 \overline{\langle f,g\rangle}=\bar f\bar g.
\]
The diagonal $\Delta_X=\langle1_X,1_X\rangle$ is total. Total arrows
have ordinary binary products. The union convention in the next section
is Heller's universal agreement convention \cite[\S6]{Heller}.

Domains $c,c'$ on $U$ are complementary if $cc'=0_U$ and
$c\vee c'=1_U$. They separate $a,b$ if $a\le c$ and $b\le c'$.
A disjoint pair is inseparable if no complementary pair separates it.
For a total negation-code map $\nu:S\to S$, a theorem domain $t$ on
$S$ is consistent if $t\nu^*t=0_S$, and complete if
$t\vee\nu^*t=1_S$.

\begin{proposition}\label{gen:turing}
Suppose a restriction category has binary restriction products and finite
domain joins preserved by total inverse image. Let $U$ satisfy
$0_U\ne1_U$, and let $\tau:U\boxtimes U\to U$ have the following
parameter-universality property: every $F:U\boxtimes U\to U$ has a
total $g:U\to U$ with $\tau(g\boxtimes1_U)=F$. Suppose disjoint
output domains $q_0,q_1$ have binary case distinction: for every
complementary $c,c'$ there is a total $h:U\to U$ with
$h^*q_0=c'$ and $h^*q_1=c$. Then, for $d=\tau\Delta_U$, the domains
$a=d^*q_0$ and $b=d^*q_1$ are inseparable.
\end{proposition}
\begin{proof}
They are disjoint because inverse image preserves meets and zero.
If $c,c'$ separated them, choose $h$ as stated and apply universality
to $h\pi_2$. Precomposition with $\langle1_U,g\rangle$ gives
$dg=hg$. Consequently
\[
 g^*a=g^*c',\qquad g^*b=g^*c.
\]
The separator inequalities make $g^*c=g^*c'$. These domains are
complementary because $g$ is total, contradicting $0_U\ne1_U$.
\end{proof}

The Turing morphism constructs this inseparable pair and, in
Section~\ref{gen:turing-fold-section}, the partial folds for finitary
inference polynomials. The finite-proof construction also uses ranges,
constructor embeddings, and domain unions. The consistency
argument additionally uses finite-list recursion and decidable equality
on the atom object.


\section{Initial algebras from domains and partial folds}\label{gen:finite-algebras}

A constructor map can encode arbitrary one-step expressions without
requiring every code to be well founded. The domain generated in finitely
many constructor steps selects the finite expressions. To obtain an
initial algebra, the corresponding finite partial folds must also assemble
to an arrow.

Let $\mathscr A$ be a restriction category with zero arrows satisfying
$\overline{0_{X,Y}}=0_X$. The restriction order is
$f\le g$ when $f=g\bar f$. A \emph{restriction functor}
$P:\mathscr A\to\mathscr A$ is a functor satisfying
$P(\bar f)=\overline{P(f)}$. It preserves the restriction order and
induces a functor
\[
 \widetilde P:\Tot(\operatorname{Split}_r\mathscr A)
 \longrightarrow\Tot(\operatorname{Split}_r\mathscr A),
 \qquad (X,x)\longmapsto(PX,Px),\quad f\longmapsto P(f).
\]
Here an object $(X,x)$ consists of a restriction idempotent on $X$,
and a total arrow $f:(X,x)\to(Y,y)$ satisfies
$f=yfx$ and $\bar f=x$.

For restriction idempotents $e_n$ on $W$, a \emph{union} $e$ means an
upper bound such that, for all parallel arrows $v,w:W\to Z$,
\[
 (\forall n,\ ve_n=we_n)\quad\Longrightarrow\quad ve=we.
\]
This condition implies that $e$ is their least upper bound. Stability
means that the same condition holds after restriction product with any
identity arrow.

Suppose there is a total arrow $\gamma:PW\to W$ and an arrow
$\delta:W\to PW$ such that
\begin{equation}\label{gen:code-inverse}
 \delta\gamma=1_{PW},\qquad
 \gamma\delta=\bar\delta.
\end{equation}
Thus $\gamma$ is a restriction monomorphism, with restriction inverse
$\delta$. Define domains on $W$ by
\begin{equation}\label{gen:finite-stages}
 e_0=0_W,\qquad e_{n+1}=\gamma P(e_n)\delta.
\end{equation}
For any restriction idempotent $v$ on $PW$,
$\gamma v\delta=\overline{v\delta}$, by the restriction identities
and \eqref{gen:code-inverse}. Hence each $e_n$ is a restriction
idempotent. They form an increasing sequence, since $P$ preserves the
restriction order. The functor $P$ need not preserve zero; its value at
$e_0$ can include nullary constructors.

\begin{theorem}\label{gen:initial}
Let $P:\mathscr A\to\mathscr A$ be a restriction functor and let
$\gamma,\delta$ satisfy \eqref{gen:code-inverse}. Assume that the
sequence \eqref{gen:finite-stages} has a union $e$ and that $P(e)$ is
a union of the domains $P(e_n)$. For every $\widetilde P$-algebra
\[
 \alpha:(PX,Px)\longrightarrow(X,x)
 \quad\text{in }\Tot(\operatorname{Split}_r\mathscr A),
\]
form the sequence of arrows $h_n:W\to X$ given by
\begin{equation}\label{gen:partial-folds}
 h_0=0_{W,X},\qquad h_{n+1}=\alpha P(h_n)\delta.
\end{equation}
Assume that this sequence has an upper bound $h:W\to X$ with
$\bar h=e$. Then
\[
 \iota=\gamma P(e):(PW,P(e))\longrightarrow(W,e)
\]
is an initial $\widetilde P$-algebra. The unique algebra morphism from
it to $\alpha$ is the arrow $h:(W,e)\to(X,x)$ obtained from
\eqref{gen:partial-folds}.
\end{theorem}
\begin{proof}
For every $n$, the equality
$e_{n+1}\gamma P(e_n)=\gamma P(e_n)$ and the inequality
$e_{n+1}\le e$ give
$e\gamma P(e_n)=\gamma P(e_n)$. The union property of $P(e)$ gives
\[
 e\gamma P(e)=\gamma P(e).
\]
Since $\gamma$ is total, $\overline{\gamma P(e)}=P(e)$.
Thus $\iota$ is a total arrow between the indicated split objects.

For a target algebra, its support and totality equations are
$x\alpha=\alpha$ and $\bar\alpha=P(x)$. Inductively,
\[
 xh_n=h_n,\qquad \bar h_n=e_n.
\]
Indeed, the first equality for $h_{n+1}$ follows from
$x\alpha=\alpha$. The inductive hypothesis gives
$P(x)P(h_n)=P(h_n)$; therefore
\[
 \bar h_{n+1}
 =\overline{P(h_n)\delta}
 =\overline{P(e_n)\delta}
 =\gamma P(e_n)\delta=e_{n+1}.
\]
The arrows $h_n$ are increasing, beginning with zero, because $P$ and
composition preserve the restriction order.

Let $h$ be the specified upper bound. Since $h_n\le h$ and
$\bar h_n=e_n$, we have $he_n=h_n$. The equalities
$xhe_n=xh_n=h_n=he_n$ and the union property of $e$ imply
$xh=xhe=he=h$. Thus $h$ is a total arrow $(W,e)\to(X,x)$.
For every $n$,
\begin{align*}
 h\gamma P(e_n)
 &=he_{n+1}\gamma P(e_n)\\
 &=h_{n+1}\gamma P(e_n)\\
 &=\alpha P(h_n)\delta\gamma P(e_n)\\
 &=\alpha P(h_n)
  =\alpha P(h)P(e_n).
\end{align*}
The union property of $P(e)$ now gives
$h\gamma P(e)=\alpha P(h)P(e)=\alpha P(h)$.
This is the algebra-morphism equation.

If $k:(W,e)\to(X,x)$ is another algebra morphism, induction gives
$ke_n=h_n$. The initial case uses zero. For the induction step,
\[
 ke_{n+1}
 =k\gamma P(e_n)\delta
 =\alpha P(k)P(e_n)\delta
 =\alpha P(ke_n)\delta=h_{n+1}.
\]
The union property gives $ke=he$, and both maps are supported on $e$.
Consequently $k=h$.
\end{proof}

The upper bound in Theorem~\ref{gen:initial} is automatically the join
of the sequence. If $v$ is another upper bound, then
$ve_n=he_n$ for every $n$, so $ve=h$ and $h\le v$.
It is enough, in particular, that the particular sequences
\eqref{gen:partial-folds} have compatible joins whose restriction is the
union of their restrictions. Preservation of arrow joins by $P$ is not
required: the proof uses the union agreement property on $P(e_n)$.

\section{Typed constructor maps from a universal object}\label{gen:coding}

Ranges turn an arbitrary retraction into a restriction inverse. This
supplies the constructor maps needed in Theorem~\ref{gen:initial}
from a suitable universal-object hypothesis.

\begin{lemma}\label{gen:section-inverse}
Suppose a range restriction category has the universal range property
\[
 vf=wf\quad\Longrightarrow\quad
 v\ran(f)=w\ran(f).
\]
Let $\gamma:Z\to W$ be total and let $r:W\to Z$ satisfy
$r\gamma=1_Z$. Put $j=\ran(\gamma)$ and $\delta=rj$. Then
\[
 \delta\gamma=1_Z,\qquad
 \gamma\delta=j=\bar\delta.
\]
\end{lemma}
\begin{proof}
The first equality follows from $j\gamma=\gamma$. Since
$\gamma r\gamma=\gamma$, the universal range property gives
$\gamma rj=j$. The restriction equation
$\bar r\gamma=\gamma\overline{r\gamma}=\gamma$ gives
$\bar rj=j$ by the same property. Hence
$\bar\delta=\overline{rj}=\bar rj=j$.
\end{proof}

Let $\mathscr A$ now be a split range restriction category with
universal ranges and binary restriction products. Put
$\mathscr E=\Tot(\mathscr A)$, and assume that it has the finite
limits and finite coproducts used below. Fix an expression object $S$.
The partial slice $\mathscr A_S$ has total arrows $q:X\to S$ as
objects. An arrow from $q:X\to S$ to $q':Y\to S$ is an arrow
$f:X\to Y$ of $\mathscr A$ satisfying
\[
 q'f=q\bar f.
\]
Restriction and composition are inherited from $\mathscr A$.
Its total subcategory is $\mathscr E/S$.

For $0\le i\le k$, supply objects $A_i$ and total maps
$s_i:A_i\to S^i$ and $t_i:A_i\to S$. These are rule instances,
premise tuples, and conclusions. Define
\[
 P_\Sigma:\mathscr E/S\longrightarrow\mathscr E/S,
 \qquad
 (q:X\to S)\longmapsto
 \left(\coprod_{i=0}^k A_i\times_{S^i}X^i
     \xrightarrow{[t_i\pi_{A_i}]}S\right).
\]
On a slice arrow it fixes the rule instance and applies that arrow to
each premise. Suppose this functor extends to a restriction functor
$P_\Sigma:\mathscr A_S\to\mathscr A_S$.

Here extensivity means that finite coproducts exist and that the canonical
functor
\[
 (\mathscr E/X)\times(\mathscr E/Y)\longrightarrow\mathscr E/(X+Y),
\]
which takes coproducts on objects and arrows, is an equivalence, with the
corresponding empty-coproduct condition. One sufficient condition for the
extension is that $\mathscr E$ be extensive and that finite products and
coproducts of total
restriction monomorphisms are again restriction monomorphisms. Let
$\mathcal M$ be their class in $\mathscr E$. It is stable under
pullback: the pullback of $m$ along a total $f$ is obtained by splitting
the domain of $m^\dagger f$. Each partial arrow is represented by its
domain span, whose left leg belongs to $\mathcal M$. Finite products,
pullback against $s_i$, and the extensive coproduct preserve these
domain spans and their composition pullbacks. Applying $P_\Sigma$ to
both legs therefore defines the required partial-arrow action. Its
restriction is obtained by applying $P_\Sigma$ to the domain inclusion,
so it is a restriction functor.

\begin{proposition}\label{gen:typed-coding}
Assume that every object of $\mathscr E$ admits a total arrow into an
object $U$ with a partial retraction. For a rule presentation $\Sigma$
as above, put $W=S\boxtimes U$ and $q=\pi_1:W\to S$.
There are arrows in $\mathscr A_S$
\[
 \gamma:P_\Sigma(q)\longrightarrow q,
 \qquad \delta:q\longrightarrow P_\Sigma(q)
\]
such that $\gamma$ is total and
$\delta\gamma=1$, $\gamma\delta=\bar\delta$.
\end{proposition}
\begin{proof}
Write $Z$ for the carrier of $P_\Sigma(q)$ and $t:Z\to S$ for its
conclusion map. Choose a total $k:Z\to U$ and a partial retraction
$\ell:U\to Z$. Define
\[
 \gamma=\langle t,k\rangle:Z\to S\boxtimes U,
 \qquad r=\ell\pi_2:S\boxtimes U\to Z.
\]
Then $q\gamma=t$ and $r\gamma=1_Z$. Lemma~\ref{gen:section-inverse}
provides $\delta=r\ran(\gamma)$ with the required inverse equations.
It remains to check its partial-slice type. Since
$t\delta\gamma=t=q\gamma$, universal ranges give
$t\delta\ran(\gamma)=q\ran(\gamma)$. The inverse equations and
$\delta=\delta\ran(\gamma)$ turn this into
$t\delta=q\bar\delta$, as required.
\end{proof}

Apply Theorem~\ref{gen:initial} in $\mathscr A_S$ to this constructor
map. If its finite-stage domains have the two indicated unions and its
partial-fold sequences admit the indicated joins, it gives an initial
$P_\Sigma$-algebra in $\mathscr E/S$. Splitting in
$\mathscr A_S$ introduces no further object requirement: a splitting
of a domain on $q:X\to S$ inherits its total map to $S$.

For a finitary polynomial, the condition on $P_\Sigma(e)$ follows
from preservation of the relevant unions by restriction products,
total inverse image, and the finite coproducts in the presentation.
Indeed, every tuple of finitely many stages is bounded by the stage
with their maximum index. Thus the union of the $i$-fold diagonal
stages is the $i$-fold product of the union. Pullback against $s_i$
and the finite coproduct give exactly the asserted union on
$P_\Sigma(e)$. The union of the finite-stage domains themselves and
the gluing of partial folds remain separate existence conditions.

\begin{corollary}\label{gen:parameter-stability}
For a typed presentation $\Sigma$ over $S$, suppose the two unions in
Theorem~\ref{gen:initial} are stable under restriction products.
Suppose also that its partial-fold gluing hypothesis holds after adjoining
any parameter object $K$. Then the constructed initial algebra is stable
under this parameter extension: if its carrier is $c:D\to S$, the
parameterized presentation over $K\boxtimes S$ has initial carrier
$1_K\boxtimes c:K\boxtimes D\to K\boxtimes S$.
\end{corollary}
\begin{proof}
The parameterized rule objects are $K\boxtimes A_i$. All premises
and the conclusion retain the same $K$-coordinate. The polynomial
\[
 P_{K\Sigma}:\mathscr E/(K\boxtimes S)
       \longrightarrow\mathscr E/(K\boxtimes S)
\]
satisfies the canonical identification
\[
 P_{K\Sigma}(1_K\boxtimes q)
 \cong 1_K\boxtimes P_\Sigma(q).
\]
Finite pullbacks and the stable coproducts give this identification on
objects and morphisms. Use the constructor $1_K\boxtimes\gamma$
and decoder $1_K\boxtimes\delta$. Their finite stages are
$1_K\boxtimes e_n$, by induction. The stable-union assumptions give
their union $1_K\boxtimes e$ and the required union of its polynomial
images. The assumed gluing for parameterized folds permits application
of Theorem~\ref{gen:initial}. Splitting $1_K\boxtimes e$ gives
$K\boxtimes D$ with its displayed conclusion map.
\end{proof}



\section{Turing morphisms and partial folds}\label{gen:turing-fold-section}

For finitary inference polynomials, a Turing morphism constructs the
partial fold required by Theorem~\ref{gen:initial}. Its restriction to
the finite-stage domain is the total fold. The argument uses an evaluator
for partial maps that respect the expression type.

Let $\mathscr A$ be a split range restriction category with restriction
zeros, binary restriction products, and universal ranges. Put
$\mathscr E=\Tot(\mathscr A)$, and assume that $\mathscr E$ has a
terminal object $1$. A total restriction embedding means a total
restriction monomorphism. Suppose every object has such an embedding
into an object $U$. Fix a Turing morphism
$\tau:U\boxtimes U\to U$ with the following parameter property:
for every object $Y$ and arrow $F:Y\boxtimes U\to U$, there is a
total $c:Y\to U$ such that
\begin{equation}\label{gen:turing-parameter}
 \tau(c\boxtimes1_U)=F.
\end{equation}

This property passes from a restriction category to its restriction
idempotent splitting. Indeed, an arrow
$F:(Y,e)\boxtimes U\to U$ has an underlying arrow
$F:Y\boxtimes U\to U$ with $F=F(e\boxtimes1_U)$.
If $c:Y\to U$ is supplied by \eqref{gen:turing-parameter}, the arrow
$ce:(Y,e)\to U$ is total in the splitting and satisfies the same
evaluation equation there. Thus a terminal object supplied by the
splitting may be used as a parameter object in that equation.

\begin{lemma}[Parameterized recursion]\label{gen:recursion}
For every arrow $G:Y\boxtimes U\boxtimes U\to U$, there is a total
$c:Y\to U$ such that
\[
 \tau\langle c\pi_1,\pi_2\rangle
 =G\langle\pi_1,c\pi_1,\pi_2\rangle:
 Y\boxtimes U\longrightarrow U.
\]
\end{lemma}
\begin{proof}
Choose a total restriction embedding
$\kappa:U\boxtimes U\to U$ with partial inverse
$\lambda:U\to U\boxtimes U$. Parameter universality supplies a
total specialization map $s:U\boxtimes U\to U$ with
\[
 \tau(s(a,b),x)=\tau(a,\kappa(b,x)).
\]
Here, and in the following calculation, expressions with arguments
denote compositions with the indicated projections and pairings.
For instance, the displayed equation is between arrows
$U\boxtimes U\boxtimes U\to U$.

Define $F:Y\boxtimes U\to U$ by
\[
 F=G(1_Y\boxtimes(s\Delta_U\boxtimes1_U))
       (1_Y\boxtimes\lambda).
\]
Choose a total $k:Y\to U$ with $\tau(k\boxtimes1_U)=F$, and set
$c=s\langle k,k\rangle:Y\to U$. For the associated arrows on
$Y\boxtimes U$, specialization and $\lambda\kappa=1$ give
\begin{align*}
 \tau(c(y),x)
 &=\tau(k(y),\kappa(k(y),x))\\
 &=F(y,\kappa(k(y),x))\\
 &=G(y,s(k(y),k(y)),x)=G(y,c(y),x).
\end{align*}
The map $c$ is total because $s$ and $k$ are total.
\end{proof}

Fix total maps $q:W\to S$ and $p:X\to S$. Choose restriction
embeddings $j_W:W\to U$ and $j_X:X\to U$, with inverses
$r_W$ and $r_X$. The arrow
\[
 H_0=r_X\tau(1_U\boxtimes j_W):U\boxtimes W\longrightarrow X
\]
evaluates program codes on $W$ and retains outputs in $X$. To restrict
it to the matching expression type, put
\[
 r=\ran(\Delta_S),\qquad
 v=\langle pH_0,q\pi_2\rangle:U\boxtimes W\to S\boxtimes S,
 \qquad e=\overline{rv},\qquad H=H_0e.
\]
The equality $\pi_1\Delta_S=\pi_2\Delta_S$ and universal ranges give
$\pi_1r=\pi_2r$. Moreover, $e\le\bar v=\bar H_0$, and the restriction
identities give $rv=ve$. Consequently
\[
 pH=\pi_1ve=\pi_1rv=\pi_2rv=\pi_2ve=q\pi_2e.
\]
Since $\bar H=e$, the arrow $H$ belongs to $\mathscr A_S$, with source
$q\pi_2$ and target $p$.

This restriction retains every well-typed specialization of $H_0$:
if $c:1\to U$ is total and $H_0(c\boxtimes1_W)=g$ with
$pg=q\bar g$, then
\begin{equation}\label{gen:type-filter}
 H(c\boxtimes1_W)=g.
\end{equation}
Indeed, write $j=c\boxtimes1_W$. The typing equation gives
$vj=\Delta_Spg$, whence $rvj=vj$ and
$j^*e=\overline{rvj}=\bar g$. Therefore
$Hj=H_0ej=H_0j\bar g=g$.

For an object $K$, write $K\cdot q$ for
$q\pi_2:K\boxtimes W\to S$. A finitary rule polynomial on
$\mathscr A_S$ has a total map
\[
 \operatorname{st}_{K,q}:K\cdot P_\Sigma(q)
                  \longrightarrow P_\Sigma(K\cdot q).
\]
It repeats the $K$-parameter at every premise of a rule instance.
At a nullary instance it forgets the parameter. On total parameter
maps it is natural; in particular, for every total $c:1\to K$,
\begin{equation}\label{gen:total-parameter-strength}
 \operatorname{st}_{K,q}(c\boxtimes1_{P_\Sigma(q)})
   =P_\Sigma(c\boxtimes1_q),
\end{equation}
with the unit isomorphisms understood. The displayed identity follows
by the premise projections, including the empty tuple for a nullary
instance. Only naturality in total parameter maps is used here.

\begin{theorem}[Turing construction of a partial fold]\label{gen:turing-fold}
Assume the categorical and Turing hypotheses of this section. Let
$P=P_\Sigma:\mathscr A_S\to\mathscr A_S$ be a finitary rule
polynomial with its restriction-functor action and the maps
\eqref{gen:total-parameter-strength}. Let
$\gamma:P(q)\to q$, $\delta:q\to P(q)$ satisfy
\eqref{gen:code-inverse}. For every total algebra
$\alpha:P(p)\to p$, there is an arrow $h:q\to p$ in
$\mathscr A_S$ such that
\begin{equation}\label{gen:fold-fixed-point}
 h=\alpha P(h)\delta.
\end{equation}
If the finite-stage domains have the two unions in
Theorem~\ref{gen:initial}, their union supports an initial algebra
in $\mathscr E/S$.
\end{theorem}
\begin{proof}
Use the typed evaluator $H:U\cdot q\to p$ just constructed, and put
\[
 G=\alpha P(H)\operatorname{st}_{U,q}
                  (1_U\boxtimes\delta):U\boxtimes W\to X.
\]
Every factor has its indicated partial-slice type, so
$pG=q\pi_2\bar G$. Form the arrow
\[
 \widehat G=j_XG(1_U\boxtimes r_W):U\boxtimes U\to U.
\]
Lemma~\ref{gen:recursion}, with parameter object $1$, supplies a
total code $c:1\to U$ satisfying
$\tau(c\boxtimes1_U)=\widehat G(c\boxtimes1_U)$.
Precomposing with $j_W$ and using $r_Wj_W=1_W$ and
$r_Xj_X=1_X$ yields
\[
 H_0(c\boxtimes1_W)=G(c\boxtimes1_W).
\]
The right side is well typed. Equation~\eqref{gen:type-filter}
therefore gives the same equality with $H$ in place of $H_0$.
Set $h=H(c\boxtimes1_W)$. By
\eqref{gen:total-parameter-strength},
\[
 h=G(c\boxtimes1_W)
   =\alpha P(H)P(c\boxtimes1_W)\delta
   =\alpha P(h)\delta.
\]

For the finite partial folds \eqref{gen:partial-folds}, induction
gives $h_n\le h$, beginning with zero and using
\eqref{gen:fold-fixed-point}. Consequently $e_n\le\bar h$ for all
$n$, and their union satisfies $e\le\bar h$. The arrow $he$ has
restriction $e$ and extends every $h_n$. It is therefore the upper
bound required by Theorem~\ref{gen:initial}.
\end{proof}

The same theorem applies to parameterized presentations over
$K\boxtimes S$. Consequently, when the two finite-stage unions are
stable under restriction products, the initial algebras satisfy
Corollary~\ref{gen:parameter-stability} without a separate gluing
assumption. The Turing construction provides a partial fold;
the union of the finite-stage domains specifies its total finite part.


\section{The uniform numerical instance}\label{gen:numerical}

Take $\mathscr A=\operatorname{Split}_r(\R_{\mathrm{rec}})$.
The category $\R_{\mathrm{rec}}$ has the positive finite powers
$\N^k$, $k\ge1$, as objects and partial recursive functions as arrows,
with extensional equality.
Objects of $\mathscr A$ may be identified with c.e. carriers, and
its arrows with partial recursive maps between them. Finite limits of
total arrows and finite tagged coproducts
exist, and every carrier has a total computable embedding into $\N$
with a partial recursive inverse. The partial-slice polynomial above
applies a partial map to every premise, being defined exactly when all
of those applications are defined.

Fix finite numerical presentations of the rule maps and of the chosen
constructor and decoder. The domains $e_n$ are uniformly c.e.: the
recurrence \eqref{gen:finite-stages} gives a procedure from an index
for $e_n$ to an index for $e_{n+1}$. Dovetailing these procedures
recognizes their union $e$. The pointwise agreement criterion verifies
that it is a stable union. Finiteness of the premise tuples gives
$P_\Sigma(e)=\bigvee_nP_\Sigma(e_n)$ by taking the maximum of their
finitely many stage indices.

For any algebra $\alpha$ in the computational slice, the arrows $h_n$
are likewise uniformly partial recursive. Dovetail their computations
on an input until one returns. Since the $h_n$ are increasing, any two
returned values agree. The resulting partial recursive map $h$ has
domain $\bigcup_n\operatorname{dom}(h_n)$, namely $e$, and extends
every $h_n$. Theorem~\ref{gen:initial} therefore constructs the
initial algebra in the numerical computational slice.

This argument allows c.e. instance carriers, with premise and conclusion
maps total on those carriers. When the instance carriers are recursive
and the constructor is the usual finite-tree encoding, the resulting
proof carrier has the decidable validity test given by recursive checking
of its nodes.

The construction is stable under a c.e. parameter carrier $K$: retain
the parameter in every intermediate tuple and in every simulation.
The same dovetailing procedure computes each parameterized fold,
uniformly in its parameter value. Thus it supplies the parameter
stability required for list, syntax, and proof algebras below.

A category is regular if it has finite limits and
regular-epimorphism--monomorphism factorizations, with regular
epimorphisms stable under pullback. A regular epimorphism is a
coequalizer of some pair.

\begin{proposition}\label{gen:numerical-regularity}
The total category of $\mathscr A$ is regular and extensive. The range
factorization of each total arrow is its regular image factorization.
Equality on $\N$ is decidable in this category.
\end{proposition}
\begin{proof}
Finite limits are computed on c.e. carriers by tuples and equality
tests; tagged finite coproducts are disjoint and stable under pullback.
Let $f:X\to Y$ be total on its c.e. source. Its image $I$ is c.e.
The surjection $X\to I$ has a computable section: enumerate $X$,
dovetail evaluations of $f$, and return a preimage when the requested
value appears. This terminates on every input in $I$. Thus $f$
factors as a split epimorphism followed by its range inclusion.
Split epimorphisms are coequalizers of their kernel pairs, and every
regular epimorphism in this factorization system is surjective and
hence split. Their pullbacks are again split epimorphisms. This proves
regularity and the image assertion. The ordinary equality test gives
the required arrow $\N^2\to1+1$.
\end{proof}


\section{A categorical construction of a consistent presentation}
\label{gen:consistency-section}

Finite Boolean valuations suffice to prove consistency of a presentation
whose positive and negative atomic axioms have disjoint supports. The
valuation used for a proof depends on its finitely many positive axiom
leaves. This permits a construction inside a computational category even
when the entire positive axiom domain has no computational complement.

Let $\R$ be a range restriction category and put
\[
 \mathscr E=\Tot(\operatorname{Split}_r(\R)).
\]
Assume that $\mathscr E$ is regular and extensive and that the
factorization of each total arrow through its restriction range is its
regular image factorization in $\mathscr E$. Thus pullback-stable regular
images agree with the computational ranges. A distinguished subobject
of $(X,e)$ is an inclusion $(X,d)\to(X,e)$ for $d\le e$.
Under the stated compatibility assumption, every subobject of
$\mathscr E$ is distinguished up to isomorphism: apply the image
factorization to a representing monomorphism. The assumption therefore
concerns a computational category, rather than a larger category with
additional noncomputational subobjects.

The domain poset of each object is consequently canonically isomorphic
to its subobject poset.
For domains $d,e$ on $X$, their join is the regular image of the
copairing of their inclusions into $X$. It is again a computational
domain. Pullbacks preserve this image and, by extensivity, the finite
coproduct used to form it. Therefore inverse image along a total arrow
preserves finite computational joins. Intersections agree with products
of restriction idempotents, and the empty join is the zero domain.

For each $X$, define the functor
\[
 L_X:\mathscr E\longrightarrow\mathscr E,
 \qquad L_X(Y)=1+X\times Y,\qquad
 L_X(f)=1_1+(1_X\times f).
\]
Assume that $\operatorname{List}(X)$ is the carrier of an initial
$L_X$-algebra, with recursion stable under an arbitrary fixed parameter
object. Write $2=1+1$, with the two
coproduct injections denoted $0,1:1\to2$. Finite extensivity defines
the Boolean operations on $2$ by their finite truth tables. Constant
Boolean arrows are denoted by their values, with the unique arrow to
$1$ suppressed.

Fix an object $U$ with decidable equality: there is an arrow
$\epsilon:U\times U\to2$ whose inverse image of $1$ is the diagonal
$\Delta_U$. Let $a,b$ be disjoint restriction domains on $U$, and let
\[
 m_A:A\longrightarrow U,\qquad m_B:B\longrightarrow U
\]
be their inclusions in the splitting. Disjointness means
$A\times_U B\cong0$.

Suppose the functor
\[
 F_U:\mathscr E\longrightarrow\mathscr E,
 \qquad F_U(X)=U+X+X^2,
\]
with its componentwise action on arrows, has an initial algebra stable
under parameters. Denote its carrier by $S$ and its constructors by
\[
 \theta:U\to S,\qquad \nu:S\to S,
 \qquad \mathbin{\Rightarrow}:S^2\to S.
\]
These constructors express atoms, negation, and implication. The
parameter-stability assumptions on initial algebras below mean that the
corresponding algebra remains initial after adjoining a fixed parameter
to every rule instance, premise, and conclusion.

Define arrows of logical axiom instances
\[
 H_1,H_3:S^2\longrightarrow S,\qquad H_2:S^3\longrightarrow S
\]
by
\begin{align*}
 H_1(\alpha,\beta)&=\alpha\Rightarrow(\beta\Rightarrow\alpha),\\
 H_2(\alpha,\beta,\gamma)
   &=(\alpha\Rightarrow(\beta\Rightarrow\gamma))
      \Rightarrow((\alpha\Rightarrow\beta)
                    \Rightarrow(\alpha\Rightarrow\gamma)),\\
 H_3(\alpha,\beta)&=(\nu\beta\Rightarrow\nu\alpha)
                         \Rightarrow(\alpha\Rightarrow\beta).
\end{align*}
The formulas here denote composites of the constructors and projections;
they do not require elements of $S$. Let the nullary rule object be
\[
 A_0=A+B+S^2+S^3+S^2,
\]
with conclusion map $[\theta m_A,\nu\theta m_B,H_1,H_2,H_3]$.
There are no unary rules. The binary rule object is $A_2=S^2$, with
premises $(\alpha,\alpha\Rightarrow\beta)$ and conclusion $\beta$.
These data determine the polynomial functor
\[
 P:\mathscr E/S\longrightarrow\mathscr E/S,
 \qquad
 P(q:X\to S)=
 \left(A_0+(A_2\times_{S^2}X^2)\longrightarrow S\right).
\]
On arrows, $P$ fixes rule instances and applies the given slice arrow
to both premise entries. Suppose this functor has a parameter-stable
initial algebra, with underlying conclusion map $c:D\to S$.

\begin{lemma}\label{gen:theorem-image}
The range $t_0=\ran(c)$ is the least distinguished subobject of $S$
containing all nullary conclusions and closed under the displayed binary
rule.
\end{lemma}
\begin{proof}
Factor $c$ as a regular epimorphism $e:D\to T_0$ followed by its range
inclusion $m:T_0\to S$. The functor $P$ preserves regular epimorphisms:
finite products and pullbacks preserve them in a regular category, and
finite coproducts preserve them in an extensive regular category. Thus
$P(e)$ is a regular epimorphism. Grafting proves that the structure
arrow of the slice object $P(m)$ factors through $m$ after
precomposition with $P(e)$.
The regular-image property therefore makes that arrow factor
through $m$ itself. This is closure under the rules. If a distinguished
subobject $n:T\to S$ is rule-closed, the factorizations of its rule
conclusions equip it with a $P$-algebra. Initiality gives a map
$D\to T$ over $S$, whence $t_0\le t$ by the leastness of the range.
\end{proof}

\begin{theorem}\label{gen:finite-valuation-consistency}
Let $\R$ be a range restriction category such that
$\mathscr E=\Tot(\operatorname{Split}_r\R)$ is regular and
extensive, restriction ranges agree with regular images, and
finite-list recursion is stable under parameters. Let $U$ have
decidable equality, and let $a,b$ be disjoint restriction domains
on $U$. Suppose that the syntax functor $F_U$ and the rule functor
$P$ defined above have parameter-stable initial algebras, with
syntax object $S$ and proof conclusion map $c:D\to S$.
Then $t_0=\ran(c)$ is consistent:
\[
 t_0\wedge\nu^*t_0=0_S.
\]
\end{theorem}
\begin{proof}
Put $L=\operatorname{List}(A)$. Decidable equality on $U$ and list
recursion define the membership arrow
\[
 \mu:U\times L\longrightarrow2.
\]
Its value is $1$ precisely when the atom agrees with the image under
$m_A$ of an entry of the list. The empty-list value is $0$, and the
cons step takes the disjunction of the equality test for the new entry
and the preceding membership value. List induction and
$A\times_U B\cong0$ give
\begin{equation}\label{gen:negative-list}
 \mu(m_B\times1_L)=0.
\end{equation}
Indeed, a successful equality test between an $A$-entry and a $B$-entry
would give an arrow to their zero pullback.

Parameterized syntax recursion gives an evaluation arrow
\[
 v:S\times L\longrightarrow2
\]
that evaluates an atom by $\mu$, negation by Boolean complementation,
and implication by its Boolean truth table. Each of $H_1,H_2,H_3$
evaluates to $1$, as is checked on the finitely many Boolean tuples.
Modus ponens preserves the value $1$.

There is also an arrow $\ell:D\to L$ listing the positive atomic
axiom leaves of a proof. At an $A$-axiom it returns the singleton list;
at all other nullary rules it returns the empty list; at a binary rule
it concatenates the lists of its premise proofs. This is a fold of the
initial proof algebra into the object $S\times L\to S$.

Membership and Boolean conjunction over finite lists define an arrow
\[
 h:D\times L\longrightarrow2
\]
which tests whether every entry of $\ell(p)$ occurs in the parameter
list. Form the decidable subobject $G\hookrightarrow D\times L$
specified by the Boolean implication
\begin{equation}\label{gen:good-proof}
 h(p,l)=1\quad\Longrightarrow\quad v(c(p),l)=1.
\end{equation}
Here and below this notation specifies identities and subobjects using
the displayed arrows; it is valid in the internal logic of
$\mathscr E$.

The object $G\to S\times L$ is a subalgebra of the proof algebra
with parameter $L$. At a positive axiom, the antecedent in
\eqref{gen:good-proof} is exactly the required atomic membership.
At a negative axiom, \eqref{gen:negative-list} makes its negation true.
At a logical axiom, Boolean validity gives the conclusion. At a binary
node, containment of the concatenated positive-leaf list implies
containment for each premise. The two premise conditions and the truth
table for implication then give the conclusion condition.

By parameter-stable initiality, the inclusion $G\to D\times L$ is
an isomorphism. Explicitly, the unique algebra map from the initial
algebra into $G$, followed by the inclusion, is its identity algebra
endomorphism. The monic inclusion is consequently a split epimorphism.
Thus \eqref{gen:good-proof} holds on all of $D\times L$.

Define the refutation object by the pullback
\[
\begin{tikzcd}
 D^-\arrow[r,"r"]\arrow[d,"c^-"']&D\arrow[d,"c"]\\
 S\arrow[r,"\nu"']&S.
\end{tikzcd}
\]
Pullback stability of the regular images gives
$\ran(c^-)=\nu^*t_0$. Let $Q=D\times_S D^-$, using $c$ and $c^-$,
with projections $p:Q\to D$ and $q:Q\to D^-$. Evaluate both
proofs against the list $\ell p\mathbin{+\!+}\ell rq$, where
$+\!+$ denotes concatenation. Each positive-leaf list is contained in
this concatenation. Hence \eqref{gen:good-proof} gives the value $1$
to both conclusions. But $crq=\nu cp$, so their evaluations are Boolean
complements. The disjoint coproduct injections $0,1:1\to2$ force
$Q\cong0$. Pullback stability of images now gives
$t_0\wedge\nu^*t_0=0_S$.
\end{proof}

\begin{corollary}\label{gen:essential}
Under the hypotheses of
Theorem~\ref{gen:finite-valuation-consistency}, assume additionally
that $a,b$ are inseparable computational domains.
Then every consistent computational domain $t$ on $S$ with
$t_0\le t$ is incomplete. In particular, this holds for every
consistent computational theory extending the displayed presentation,
where a computational theory means a distinguished subobject closed
under its logical axioms and modus ponens.
\end{corollary}
\begin{proof}
The nullary constructors give arrows $u_A:A\to D$ and $u_B:B\to D$
with
\[
 cu_A=\theta m_A,\qquad cu_B=\nu\theta m_B.
\]
Taking ranges gives
$a\le\theta^*t_0$ and $b\le\theta^*\nu^*t_0$.
If an extension $t$ were consistent and complete, then
$t\wedge\nu^*t=0_S$ and $t\vee\nu^*t=1_S$. Total inverse image
along $\theta$ would give a complementary separator of $a,b$.
\end{proof}



\section{Presheaf completion and the range obstruction}
\label{gen:completion}

Heller constructs a recursion category from partial maps of simplicial
sets in which the diagonal $Y\to Y\boxtimes Y$ has no range
\cite[\S10, Example~5, p.~1267]{Heller}. Range existence must therefore
be established in the category used for deduction. The following
calculation identifies a failure in a selected presheaf completion.

Let $\mathscr C$ be a small
category with a stable class $\mathcal M$ of monomorphisms: it contains
isomorphisms and is closed under composition and pullback. Write
$y:\mathscr C\to\Set^{\mathscr C^{\op}}$ for Yoneda. Select a
presheaf monomorphism $\mu:P\to Q$ when each pullback along
$yX\to Q$ is isomorphic to $ym:yA\to yX$ for some $m\in\mathcal M$.
Denote this class by $\widehat{\mathcal M}$.
Garner and Lin prove that it is stable and that Yoneda induces
\begin{equation}\label{gen:selected-domains}
 \operatorname{Sub}_{\mathcal M}(X)
 \cong\operatorname{Sub}_{\widehat{\mathcal M}}(yX)
\end{equation}
\cite[\S3.1 and Lemma~9]{GarnerLin}.
Applying the defining pullback condition to $1_{yX}$ gives surjectivity;
full faithfulness gives the order isomorphism. Thus equations for
complementary separators in these selected orders are preserved and
reflected. If $0$ is initial in $\mathscr C$ and every initial map
$0\to X$ belongs to $\mathcal M$, the selected bottom on
$yX$ is $y0\to yX$; $y0$ need not be the empty presheaf.

Write $\operatorname{Par}(\mathscr E,\mathcal M)$ for the category
whose arrows $X\rightharpoonup Y$ are spans
$X\xleftarrow{m}D\xrightarrow{f}Y$ with $m\in\mathcal M$, modulo
isomorphism of spans. Composition uses pullbacks, and restriction
replaces the right leg $f$ by $m$. Its total arrows are ordinary arrows
of $\mathscr E$. The following obstruction concerns the existence of
ranges in this partial-map category.

\begin{proposition}\label{gen:no-presheaf-range}
Let $\mathscr C_{\mathrm{rec}}$ be the total category of the numerical
restriction-idempotent splitting, and let $\mathcal M$ consist of its
restriction monomorphisms. For a non-c.e. set $B\subseteq\N$, the
copairing of the numeral maps
\[
 F_B:\coprod_{n\in B}y1\longrightarrow y\N
\]
has no restriction range in
$\operatorname{Par}(\Set^{\mathscr C_{\mathrm{rec}}^{\op}},
\widehat{\mathcal M})$.
\end{proposition}
\begin{proof}
A range would, by \eqref{gen:selected-domains}, correspond to the
inclusion of a c.e. set $T\subseteq\N$. Every summand factors through
it, so $B\subseteq T$ by full faithfulness of Yoneda. For $k\notin B$,
each numeral summand factors through $y(\N\setminus\{k\})$. The
coproduct property gives the same factorization for $F_B$. Leastness of
a range therefore gives $T\subseteq\N\setminus\{k\}$ for every
$k\notin B$. Hence $T=B$, a contradiction.
\end{proof}

This obstruction uses an externally indexed coproduct. The finite-stage
construction of Theorem~\ref{gen:initial} requires unions for specified
recursively generated domain sequences. In the numerical case their
uniformity makes these unions available. Theorem~\ref{gen:turing-fold}
supplies the partial folds without a separate arrow-gluing assumption.
It does not supply the finite-stage domain unions.

\begin{problem}\label{gen:next-problem}
For a recursion category satisfying the hypotheses of
Theorem~\ref{gen:turing-fold}, derive the required finite-stage domain
unions from a categorical iteration structure. Construct a finite-tree
evaluator and prove that its termination domain is the union of the
constructor stages, with the stated stability and polynomial preservation.
Determine which such categories satisfy the regular-image and finite
valuation hypotheses of Theorem~\ref{gen:finite-valuation-consistency}.
\end{problem}

The fixed-point equation \eqref{gen:fold-fixed-point} determines the
fold on every finite stage, but it does not assert that the fold's full
domain is generated by those stages. The domain equality required in
this problem is therefore additional to the recursion theorem.

\begin{thebibliography}{9}
\bibitem{CockettLack} J.~R.~B.~Cockett and S.~Lack,
Restriction categories I: categories of partial maps,
\emph{Theoretical Computer Science} \textbf{270}(1--2) (2002), 223--259.
\bibitem{Heller} A.~Heller,
An existence theorem for recursion categories,
\emph{Journal of Symbolic Logic} \textbf{55}(3) (1990), 1252--1268.
\href{https://doi.org/10.2307/2274486}{doi:10.2307/2274486}.
\bibitem{GarnerLin} R.~Garner and D.~Lin,
\emph{Cocompletion of restriction categories},
\href{https://arxiv.org/abs/1610.07164}{arXiv:1610.07164}.
\end{thebibliography}
\end{document}
